\chapter{语法知识}


\section*{词类}
词类有实词与虚词两大类。

\subsection*{实词}
表示实在的意义，能够作短语或句子的成分，能够独立成句。

\subsection*{虚词}
一般不表示实在的意义，不作短语或句子的成分（只有副词例外），它们的基本用途是表示语法关系。

\section{实词分类}

\subsection{A、名词}
表示人和事物名称的词。

\begin{itemize}
    \item 表示人的名称，如：同志、作家
    \item 表示具体事物，如：河流、高山
    \item 表示抽象事物，如：政治、科学
    \item 表示时间名称，如：上午、夏天
    \item 表示处所名称，如：上海、中国
    \item 表示方位名称，如：上、下（简称方位词）
\end{itemize}

\subsubsection*{名词的语法特点}
\begin{enumerate}
    \item 表示人称的名词，可以在后头加“们”表示多数。
    \item 方位词常用在其他名词后头，组成表示处所、范围或时间的方位短语。
    \item 名词一般不受副词修饰。
\end{enumerate}

\subsection{B、动词}
表示动作行为、发展变化、心理活动等意义的词。

\begin{itemize}
    \item 表示动作、行为，如：坐、听
    \item 表示存现、消失或发展变化，如：有、发生
    \item 表示心理活动，如：爱、恨
    \item 表示使令，如：叫、让
    \item 表示可能、意愿（能愿动词），如：能、会
    \item 表示趋向（趋向动词），如：来、去
    \item 表示判断（判断词），如：是
\end{itemize}

\subsubsection*{动词的语法特点}
\begin{enumerate}
    \item 动词一般受副词“不”的修饰。
    \item 动词后面可以带“着、了、过”，表示动态。
    \item 一部分动词可以重叠，表示时间短暂或尝试的意思。
    \item 判断词“是”主要是联结句子的主语和宾语。
    \item 能愿动词后面不能跟名词，能愿动词可以和后面的动词一起作谓语中心语，也可以单独作谓语中心语。
    \item 趋向动词可以单独作谓语中心语，也可以在谓语中心语后面作补语。
\end{enumerate}

\subsection{C、形容词}
表示事物的形状、性质、状态的词。

\begin{itemize}
    \item 表示事物形状的，如：高、矮
    \item 表示事物性质的，如：漂亮、结实
    \item 表示事物状态的，如：快、慢
\end{itemize}

\subsubsection*{形容词的语法特点}
\begin{enumerate}
    \item 一部分形容词可以用重叠形式来加强语义。
    \item 大多数形容词可以受副词“很”修饰。
\end{enumerate}

\subsection{D、数词}
表示数目的词。

\begin{itemize}
    \item 表确数（表示分数、整数和倍数）
    \item 表概数，如：几、许多
    \item 表序数，如：第一、老三
\end{itemize}

\subsubsection*{数词的语法特点}
\begin{enumerate}
    \item 数目增加，可以用分数表示，也可以用倍数表示。
    \item 数目减少，只能用分数，不能用倍数。
\end{enumerate}

\subsection{E、量词}
表示事物和动作、行为单位的词。

\begin{itemize}
    \item 表示事物单位的量词叫物量词，如：个、只
    \item 表示动作、行为单位的量词叫动量词，如：次、回
    \item 有时也借用某些名词来表示，如：脚、年
\end{itemize}

\subsubsection*{量词的语法特点}
\begin{enumerate}
    \item 量词经常和数词连用，组成数量短语。
    \item 表示物量的数量短语常用在名词的前面。
    \item 表示动量的数量短语常用在动词的后面。
\end{enumerate}

\subsection{F、代词}
起代替或指示作用的词。

代词分为人称代词、疑问代词、指示代词三类。

\subsubsection*{代词的语法特点}
\begin{enumerate}
    \item 第二人称的敬称“您”不用于复数，如果需要表示复数，就用“您几位”“您诸位”。
    \item 第三人称复数代词“他们”可专指男性，也可兼指男性和女性，“她们”则专指女性。
    \item 注意“我们”和“咱们”用法的区别。“我们”指说话人，有时也可以包括听话人；“咱们”一定包括说话人和听话人。
    \item 指示代词“那”用于远指，“这”用于近指。
    \item 代词用得不恰当，指代不明，可造成病句。
\end{enumerate}

\section{虚词分类}

\subsection{G、副词}
一般用在动词、形容词前边，表示行为、动作或性质、状态的程度、范围、时间、频率、情势、语气等。

\begin{itemize}
    \item 表示范围，如：都、全
    \item 表示语气，如：可、倒
    \item 表示否定，如：不、没
    \item 表示时间，如：刚、恰好
    \item 表示程度，如：很、极
    \item 表示情势，如：仿佛、渐渐
\end{itemize}

\subsubsection*{副词的语法特点}
\begin{enumerate}
    \item 副词主要用来修饰、限制动词或形容词，在动词、形容词前面作状语。
    \item 副词有时用在形容词后面，补充说明程度、结果，作补语。
    \item 副词不能修饰名词、代词。
\end{enumerate}

\subsection{H、连词}
用来连接词、短语或句子的词。

\begin{itemize}
    \item 一般连词，如：和、与、并、或、及
    \item 关联词，如：不但……而且、虽然……但是
\end{itemize}

\subsubsection*{连词的语法特点}
\begin{enumerate}
    \item 一般连词的前后两部分可以调换而基本意思不变。
    \item 关联词主要在复句中进行运用。
\end{enumerate}

\subsection{I、介词}
经常用在名词、代词等的前面，和这些词合起来，表示动作、行为、性状的起止、方向、处所、时间、对象、方式、原因、目的、比较等。

\subsubsection*{常用介词及其用法（顺口溜）}
\begin{center}
    自、从、以、当、为、按照，\\
    由于、对于、为了、到，\\
    和、跟、把、比、在、关于，\\
    除了、同、对、向、往、朝……
\end{center}

用在名词、代词前，组成介宾短语后，修饰、补充“动”“形”要记牢。

\subsection{J、助词}
附着在实词、短语或句子上面，起辅助作用的词。

\begin{itemize}
    \item 结构助词：的、地、得
    \item 动态助词：着、了、过
    \item 语气助词：吗、吧、呢等
\end{itemize}

\subsection{K、叹词}
表示感叹、呼唤、应答等声音的词，如：啊、嗯等。

\subsubsection*{语法特点}
一般独立成句，用逗号或感叹号隔开。

\subsection{L、拟声词}
摹拟人或事物的声音的词。

\subsubsection*{语法特点}
在句子中相当于一个形容词。

\section{词类的辨别}
\begin{enumerate}
    \item \textbf{区分名词和非名词：} 名词前不能加“不”和“很”。
    \item \textbf{区分形容词和动词：} 形容词可以用“很”来修饰，动词前不能加“很”（表示心理活动的动词除外）。
    \item \textbf{区分形容词和副词：} 形容词能修饰名词，前面能加“很”；副词不能修饰名词，前面不能加“很”。
    \item \textbf{区分连词和介词：} 前后能互换的是连词，前后不能互换的是介词。
    \item \textbf{区分动词和介词：} 作谓语中心语的只能是动词，组成介宾短语修饰、补充动词、形容词的是介词。
    \item \textbf{区分语气助词和叹词：} 语气助词一般用在句尾，叹词往往独立成句，一般在句首。
    \item \textbf{区分介词和副词：} 介词后面跟名词、代词，副词后面是动词或形容词。
\end{enumerate}

\section{短语}
由词和词组合而成的语言单位。

\subsection{（1）并列短语}
由两个或由两个以上的名词、动词或形容词等并列组成的短语。

基本结构：名+名、名+代、代+代、动+动、形+形、数量+数量。

\subsubsection*{特点}
\begin{enumerate}
    \item 并列短语前后的词性一致。（名词和代词除外）
    \item 并列短语两部分之间是平等关系，没有修饰、限制关系。
    \item 并列短语中的词一般颠倒过来意思不变。
    \item 并列短语中词和词之间可以直接组合，也可以借用虚词组合。
\end{enumerate}

\subsection{（2）偏正短语}
\subsubsection*{基本结构}
\begin{itemize}
    \item 中心语是名词时，修饰限制成分是定语，用（ ）表示。结构如：形+名、数量+名、名+名、代+名。
    \item 中心语是动词或形容词时，修饰语是状语，用〔 〕表示。结构如：形+动、副+动、数量+动、副+形。
\end{itemize}

\subsection{（3）动宾短语}
动词后边带上一个受动词支配的词，组成一个短语。

基本结构：动+名、动+代。

\subsubsection*{特点}
\begin{enumerate}
    \item 动宾短语前边的动词直接支配后边的名词、代词，后边的名词、代词受前边的动词的支配，它们之间是支配和被支配的关系。
    \item 动宾短语中受动词支配的名词、代词，是宾语。
    \item 宾语一般在动词后面回答“谁”、“什么”的问题。
    \item 使用动词短语时，要注意动词和宾语意义上的配合，否则造成动宾不搭配。
\end{enumerate}

\subsection{（4）补充短语}
包括动补短语和形补短语两大类。

\subsubsection*{语法特点}
\begin{enumerate}
    \item 在动词、形容词后面起补充、说明作用的成分是补语，用＜ ＞表示。
    \item 这类短语的中心语在前，前后两部分是被补充和补充的关系。
    \item 补语在动词或形容词后边补充说明怎么样、多久、多少等问题。
    \item 有的补语前头常用结构助词“得”。
\end{enumerate}

\subsection{（5）主谓短语}
基本结构：名（代）+动、名（代）+形、名（代）+疑问代词。

特殊情况：名+名（如：今天星期一）、名+数量（如：纸三张）。

\subsubsection*{特点}
\begin{enumerate}
    \item 主谓短语前边的词表示“谁”或“什么”，后面的词说明前边的词“怎么样”“干什么”或“是什么”。前后两部分是被陈述和陈述的关系。
    \item 使用主谓短语加上语气，书面上加上标点就是一个单句，表达的意思是完整的。
\end{enumerate}

\subsection{（6）介宾短语}
由介词和它的宾语构成的短语。

基本结构：介词+名词、介词+代词。

\subsubsection*{语法特点}
\begin{enumerate}
    \item 介宾短语在句子中作为一个整体充当句子成分。
    \item 介宾短语在谓语中心语前做状语，在谓语中心语后面作补语。
    \item 介宾短语有时也做定语，后头必须带“的”。
\end{enumerate}

\subsection{（7）“的”字短语}
由动词、形容词、动宾短语加上“的”构成。

基本结构：动词+的、形容词+的、动宾短语+的。

\subsubsection*{特点}
\begin{enumerate}
    \item “的”字短语在句中相当于一个名词。
    \item “的”字短语一般常做主语、宾语。
\end{enumerate}